\documentclass{article}
\usepackage[T1]{fontenc}
\usepackage{lmodern}
\usepackage{graphicx}
\usepackage{amsmath,amssymb}
\usepackage[a4paper,margin=2cm]{geometry}
\usepackage[absolute,overlay]{textpos}
\setlength{\TPHorizModule}{1mm}
\setlength{\TPVertModule}{1mm}
\usepackage[indonesian]{babel}
\usepackage{hyperref}
\setlength{\emergencystretch}{2em}
% unit: Halaman 24

\begin{document}

% === Halaman 24 (python/native) ===

Kurva Eliptik pada Galois Field

• Operasi kurva eliptik yang dibahas sebelum ini didefinisikan pada bilangan riil.

• Operasi pada bilangan riil tidak akurat karena mengandung pembulatan

• Pada sisi lain, kriptografi dioperasikan pada ranah bilangan integer.

• Agar kurva eliptik dapat dipakai di dalam kriptografi, maka kurva eliptik 

didefinisikan pada medan berhingga atau Galois Field, yaitu GF(p) dan GF(2m).

• Yang dibahas dalam kuliah ini hanya kurva eliptik pada GF(p)

24

\end{document}
